%% ============================================================
%%  AI Driven Intelligent Rubber Production Monitoring System
%%
%%  IEEE Conference Paper — IEEEtran Document Class
%%
%%  HOW TO COMPILE:
%%    Recommended : Upload to https://www.overleaf.com
%%                  IEEEtran.cls is pre-installed there.
%%    Local       : Requires texlive-publishers package.
%%                  Run: pdflatex -> pdflatex (twice)
%% ============================================================
\documentclass[10pt, conference]{IEEEtran}
\usepackage{etoolbox}
\usepackage[T1]{fontenc}
\usepackage[utf8]{inputenc}
\usepackage{cite}
\usepackage{amsmath,amssymb,amsfonts}
\usepackage{graphicx}
\usepackage{textcomp}
\usepackage{xcolor}
\usepackage{booktabs}
\usepackage{multirow}
\usepackage{array}
\usepackage{url}
\usepackage{float}
\usepackage{microtype}
\usepackage[hidelinks]{hyperref}
\usepackage{subcaption}
\graphicspath{{figures/}}

\def\BibTeX{{\rm B\kern-.05em{\sc i\kern-.025em b}\kern-.08em
    T\kern-.1667em\lower.7ex\hbox{E}\kern-.125emX}}

\begin{document}

% ── Title ────────────────────────────────────────────────────
\title{AI Driven Intelligent Rubber Production\\Monitoring System}

% ── Authors — 3 columns × 2 rows ──────────────────────────────────────────────
\author{

% ── Row 1, Column 1 ───────────────────────────────────────────────────────────
\IEEEauthorblockN{1\textsuperscript{st} Bhanuka K.Y.}
\IEEEauthorblockA{
  \textit{dept. of Information Technology} \\
  \textit{Sri Lanka Institute of Information Technology} \\
  Malabe, Sri Lanka \\
  IT22563750@my.sliit.lk}

\and

% ── Row 1, Column 2 ───────────────────────────────────────────────────────────
\IEEEauthorblockN{2\textsuperscript{nd} Hewamanna K.T.}
\IEEEauthorblockA{
  \textit{dept. of Information Technology} \\
  \textit{Sri Lanka Institute of Information Technology} \\
  Malabe, Sri Lanka \\
  IT22584090@my.sliit.lk}

\and

% ── Row 1, Column 3 ───────────────────────────────────────────────────────────
\IEEEauthorblockN{3\textsuperscript{rd} Wadanambi P.W.}
\IEEEauthorblockA{
  \textit{dept. of Information Technology} \\
  \textit{Sri Lanka Institute of Information Technology} \\
  Malabe, Sri Lanka \\
  ITXXXXXXXX@my.sliit.lk}

\and

% ── Row 2, Column 1 ───────────────────────────────────────────────────────────
\IEEEauthorblockN{4\textsuperscript{th} Appuhamy M.H.M.M.}
\IEEEauthorblockA{
  \textit{dept. of Information Technology} \\
  \textit{Sri Lanka Institute of Information Technology} \\
  Malabe, Sri Lanka \\
  ITXXXXXXXX@my.sliit.lk}

\and

% ── Row 2, Column 2 ───────────────────────────────────────────────────────────
\IEEEauthorblockN{5\textsuperscript{th} Mr. Ravi Supunya}
\IEEEauthorblockA{
  \textit{dept. of Information Technology} \\
  \textit{Sri Lanka Institute of Information Technology} \\
  Malabe, Sri Lanka \\
  ravi.s@sliit.lk}

\and

% ── Row 2, Column 3 ───────────────────────────────────────────────────────────
\IEEEauthorblockN{6\textsuperscript{th} Mrs. Dulani Jayasinghe}
\IEEEauthorblockA{
  \textit{dept. of Information Technology} \\
  \textit{Sri Lanka Institute of Information Technology} \\
  Malabe, Sri Lanka \\
  dulani.j@sliit.lk}

}

\maketitle

% ── Abstract ─────────────────────────────────────────────────
\begin{abstract}
Sri Lanka's rubber industry, which sustains over 200,000 smallholder
families, faces persistent challenges including yield losses from
diseases and pests, subjective manual quality grading, opaque pricing,
and the absence of supply chain traceability. This paper presents the
design, implementation, and evaluation of an AI-driven intelligent
rubber production monitoring system that unifies six integrated
modules within a single mobile platform: (i)~deep-learning-based
disease, pest, and weed detection using ONNX-optimised convolutional
neural network models with GPS-tagged geospatial outbreak alerting;
(ii)~automated rubber sheet quality grading into RSS standards via
image classification; (iii)~machine-learning-based rubber price
forecasting with a blockchain-backed real-time auction marketplace;
(iv)~a Digital Product Passport (DPP) employing AES-256-CBC
zero-knowledge encryption and QR-code anti-fraud verification for
end-to-end supply chain traceability; (v)~IoT-based environmental
stress monitoring using ESP32 sensors for soil moisture, humidity, and
temperature; and (vi)~an NLP-powered agricultural chatbot for rubber
cultivation advisory. The system is built on React Native (Expo) for
cross-platform mobile delivery, ASP.NET Core~8.0 with MongoDB for the
backend API, and ONNX Runtime for on-device and server-side inference.
Experimental evaluation demonstrates 91.4\% disease classification
accuracy, 89.7\% grading accuracy, a price prediction MAE of LKR~18.3
per kilogram, and a System Usability Scale score of 79.2/100 from
field validation with 40~rubber industry stakeholders across three
districts. The system provides a scalable, technically validated
solution for modernising rubber production management in Sri Lanka.
\end{abstract}

\begin{IEEEkeywords}
rubber production, deep learning, ONNX, disease detection, rubber
grading, price forecasting, blockchain, digital product passport,
IoT, ESP32, supply chain traceability, zero-knowledge encryption,
mobile application, precision agriculture, Sri Lanka.
\end{IEEEkeywords}

% =============================================================
\section{Introduction}
% =============================================================

Rubber (\textit{Hevea brasiliensis}) is one of the most economically
significant plantation crops in Sri Lanka, supporting the livelihoods
of over 200,000 smallholder farming families and contributing
substantially to national export earnings~\cite{cbsl2024}. Despite its
importance, the industry faces systemic inefficiencies at every stage
of the production value chain---from plantation health monitoring
through quality assessment to market pricing and product
traceability.

At the plantation level, biotic threats including leaf diseases (white
root disease, anthracnose, powdery mildew, leaf spot), pests (scale
insects, mites, caterpillars), and weed competition cause yield losses
estimated at 30--40\%~\cite{rrisl2024}. Conventional identification
relies on expensive expert consultation that is largely inaccessible
to remote rural farmers~\cite{doa2024}. Post-harvest, rubber sheet
quality grading into Ribbed Smoked Sheet (RSS) categories remains a
manual, subjective process susceptible to inter-grader variability.
Price negotiation lacks transparency, and the absence of standardised
traceability mechanisms limits Sri Lankan rubber's competitiveness in
quality-conscious international markets~\cite{jayasinghe2024}.

Recent advances in deep learning, edge computing, Internet of Things
(IoT), and blockchain technology present an opportunity to address
these challenges through an integrated intelligent monitoring
platform. Lightweight CNN architectures optimised for mobile
inference, combined with the Open Neural Network Exchange (ONNX)
runtime, now enable real-time AI classification on commodity
hardware~\cite{tan2019}. Blockchain-based product passports and
zero-knowledge encryption can establish immutable trust mechanisms for
supply chain integrity~\cite{casino2019}.

This paper presents the design, implementation, and evaluation of an
\textbf{AI Driven Intelligent Rubber Production Monitoring System}
that integrates six modules into a unified mobile platform:

\begin{enumerate}
  \item \textbf{Disease, Pest, and Weed Detection} using ONNX CNN
        models with GPS geospatial outbreak alerting.
  \item \textbf{Rubber Sheet Quality Grading} via image-based
        classification into RSS~1/2/3 standards.
  \item \textbf{ML-Based Price Forecasting} with a blockchain-backed
        real-time auction marketplace.
  \item \textbf{Digital Product Passport (DPP)} with AES-256
        zero-knowledge encryption and QR anti-fraud verification.
  \item \textbf{IoT Environmental Monitoring} via ESP32 sensors for
        soil moisture, humidity, and temperature.
  \item \textbf{AI Agricultural Chatbot} for rubber cultivation
        advisory support.
\end{enumerate}

The remainder of this paper is organised as follows: Section~II
reviews related work. Section~III describes the system architecture
and methodology. Section~IV details implementation. Section~V presents
results and evaluation. Section~VI discusses findings and limitations.
Section~VII concludes the paper.

% =============================================================
\section{Related Work}
% =============================================================

\subsection{AI-Based Plant Disease Detection}

Deep learning for plant disease identification has advanced
significantly since Mohanty~\textit{et~al.}~\cite{mohanty2016}
demonstrated CNN accuracy exceeding 99\% on controlled PlantVillage
images. Withana~\cite{withana2023} developed a CNN-based rubber
disease detection system for Sri Lankan conditions achieving 87\%
accuracy, but was limited to disease classes only.
Wang~\textit{et~al.}~\cite{wang2022} proposed GMA-Net for rubber leaf
disease recognition at 92\% accuracy, emphasising region-specific
training data requirements.

\subsection{Agricultural Quality Assessment}

Automated crop quality grading using computer vision has been explored
for tea~\cite{chen2021}, rice~\cite{fabiyi2020}, and rubber latex
properties~\cite{nsbm2023}. However, image-based rubber sheet grading
into RSS categories has received limited attention, with most systems
relying on manual sensory evaluation.

\subsection{Price Prediction and Blockchain in Agriculture}

Machine learning regression models for agricultural commodity pricing
have been applied to rubber~\cite{li2022} and palm oil markets.
Blockchain-based supply chain traceability has been explored for
food~\cite{casino2019} and agricultural products, though rubber
industry adoption remains nascent. Kim~\textit{et~al.}~\cite{kim2020}
demonstrated NFT-based product passports for supply chain authenticity
verification.

\subsection{IoT in Plantation Monitoring}

Rathnayake~\textit{et~al.}~\cite{rathnayake2021} developed an
IoT-based smart agriculture monitoring system for Sri Lanka using
sensor networks, while precision agriculture platforms integrating
environmental sensing with AI advisory have shown promise in
large-scale crop management~\cite{wolfert2017}.

\subsection{Research Gaps}

Key gaps identified include: (G1)~no unified platform integrating
disease detection, quality grading, price forecasting, and
traceability; (G2)~absence of zero-knowledge encryption for rubber
supply chain documents; (G3)~lack of blockchain-backed rubber auction
mechanisms; (G4)~limited IoT integration for real-time plantation
monitoring within AI-powered rubber management systems.

% =============================================================
\section{System Architecture and Methodology}
% =============================================================

\subsection{Overall Architecture}

The proposed system comprises a four-layer architecture as illustrated
in Fig.~\ref{fig:architecture}:
\textit{Layer~1---Mobile Application:} React Native (Expo)
cross-platform client providing all user-facing modules.
\textit{Layer~2---API Gateway:} ASP.NET Core~8.0 RESTful backend with
JWT authentication, role-based access control (Farmer, Buyer,
Exporter, Admin, Researcher), and SignalR WebSocket for real-time
communication.
\textit{Layer~3---AI/ML Service:} ONNX Runtime inference engine
hosting four specialised models (disease, pest, grading, price).
\textit{Layer~4---Data Layer:} MongoDB document store for
unstructured data, user profiles, and transaction records.

%% ── FIGURE 1: System Architecture Diagram ──────────────────────
%% TODO: Replace 'system_architecture.png' with your actual diagram.
%%       Place the image file inside the 'figures/' folder.
\begin{figure*}[!t]
\centering
\includegraphics[width=\textwidth]{system_architecture.png}
\caption{Overall system architecture of the AI Driven Intelligent Rubber
Production Monitoring System, showing the four-layer design: Mobile
Application (React Native), API Gateway (ASP.NET Core~8.0), AI/ML
Service (ONNX Runtime), and Data Layer (MongoDB).}
\label{fig:architecture}
\end{figure*}

\subsection{Disease, Pest, and Weed Detection Module}

\subsubsection{ONNX Classification Pipeline}
Three independent ONNX models handle leaf disease classification
(5~classes: white root disease, anthracnose, powdery mildew, leaf
spot, healthy), pest detection (4~classes: scale insects, mites,
caterpillars, none), and weed identification (3~classes: broadleaf,
grass, sedge). Each model accepts 224${\times}$224-pixel RGB input
and produces softmax probability distributions. Predictions below
configurable confidence thresholds (0.65--0.75) are flagged as
``Unidentified'' to prevent unreliable classifications.

\subsubsection{GPS Geospatial Alert System}
Every confirmed detection is geo-tagged with device GPS coordinates
and submitted to the backend. A proximity-based alert engine computes
distances between new detections and registered farmer locations,
dispatching push notifications to farmers within configurable radii.
A real-time disease map visualises all detections with colour-coded
markers (red for leaf diseases, pink for pests, orange for weeds) and
a 5~km danger zone circle around each farmer's plantation.
Fig.~\ref{fig:disease} shows the disease detection user interface.

%% ── FIGURE 2: Disease Detection Screenshots ───────────────────
%% TODO: Replace image filenames with your actual screenshots.
%%       Place all image files inside the 'figures/' folder.
\begin{figure}[!t]
\centering
\begin{subfigure}[b]{0.48\columnwidth}
  \centering
  \includegraphics[width=\textwidth]{disease_home.png}
  \caption{Detection home}
\end{subfigure}
\hfill
\begin{subfigure}[b]{0.48\columnwidth}
  \centering
  \includegraphics[width=\textwidth]{disease_result.png}
  \caption{Detection result}
\end{subfigure}
\vskip 4pt
\begin{subfigure}[b]{0.48\columnwidth}
  \centering
  \includegraphics[width=\textwidth]{disease_map.png}
  \caption{Geospatial disease map}
\end{subfigure}
\hfill
\begin{subfigure}[b]{0.48\columnwidth}
  \centering
  \includegraphics[width=\textwidth]{disease_alerts.png}
  \caption{Proximity alerts}
\end{subfigure}
\caption{Disease, pest, and weed detection module: (a)~home screen with
detection type selection, (b)~AI classification result with severity and
remedy, (c)~geospatial map with colour-coded detection markers and 5~km
danger zone, (d)~proximity-based outbreak alert notifications.}
\label{fig:disease}
\end{figure}

\subsection{Rubber Sheet Quality Grading Module}

An ONNX image classification model grades rubber sheet photographs
into RSS~1 (Good quality), RSS~2 (Pin-hole defects), and RSS~3
(Reaper-cut defects) categories. The model outputs a predicted class,
confidence score, severity indicator, and textual improvement
suggestions. The grading screen (Fig.~\ref{fig:grading}) captures
batch metadata (tester name, batch ID, sheet count, weight) and
generates downloadable PDF quality reports for record-keeping and
compliance.

%% ── FIGURE 3: Rubber Sheet Grading UI ─────────────────────────
%% TODO: Replace image filenames with your actual screenshots.
\begin{figure}[!t]
\centering
\begin{subfigure}[b]{0.48\columnwidth}
  \centering
  \includegraphics[width=\textwidth]{grading_input.png}
  \caption{Batch info \& image capture}
\end{subfigure}
\hfill
\begin{subfigure}[b]{0.48\columnwidth}
  \centering
  \includegraphics[width=\textwidth]{grading_result.png}
  \caption{Grading result \& report}
\end{subfigure}
\caption{Rubber sheet grading module: (a)~batch information form with
image capture for quality analysis, (b)~AI grading result showing RSS
grade, confidence score, severity, and PDF report generation.}
\label{fig:grading}
\end{figure}

\subsection{ML-Based Price Forecasting Module}

\subsubsection{Price Prediction Model}
An ONNX regression model predicts per-kilogram rubber sheet prices
based on seven input features: rubber sheet grade (RSS1--RSS5),
quantity (kg), moisture level (Dry/Normal/Wet), cleanliness
(Clean/Slight/Dirty), visual quality score (1--5), production
district, and market availability timeline. The model outputs
predicted price in LKR.

\subsubsection{Blockchain-Based Auction Platform}
A real-time auction mechanism enables farmers to list rubber lots for
competitive bidding. Auctions operate on a strict 1-hour window with
SignalR WebSockets providing live bid updates across all connected
clients (Fig.~\ref{fig:price}). Each lot is tokenised as an NFT with
metadata (grade, weight, location, moisture, cleanliness) stored on
IPFS for immutable provenance. Role-based restrictions prevent farmers
from bidding on their own lots to deter price manipulation.

%% ── FIGURE 4: Price Forecasting & Auction UI ──────────────────
%% TODO: Replace image filenames with your actual screenshots.
\begin{figure}[!t]
\centering
\begin{subfigure}[b]{0.48\columnwidth}
  \centering
  \includegraphics[width=\textwidth]{price_forecast.png}
  \caption{Price calculator}
\end{subfigure}
\hfill
\begin{subfigure}[b]{0.48\columnwidth}
  \centering
  \includegraphics[width=\textwidth]{auction_bidding.png}
  \caption{Live auction bidding}
\end{subfigure}
\caption{Price forecasting and auction module: (a)~ML-based price
calculator with grade, quantity, and quality inputs, (b)~real-time
auction bidding screen with live SignalR updates, NFT passport
verification, and countdown timer.}
\label{fig:price}
\end{figure}

\subsection{Digital Product Passport Module}

\subsubsection{Zero-Knowledge Encryption Model}
Confidential trade documents are classified by an ONNX document
classifier and encrypted using AES-256-CBC with keys derived via
PBKDF2-SHA256 (100,000 iterations). The encryption key
(\textit{SecretRequestId}) is a 32-byte cryptographically secure hex
string delivered to authorised parties via Relationship-Based Access
Control (ReBAC). Client-side decryption using CryptoJS ensures
plaintext never traverses the network.

\subsubsection{QR Anti-Fraud Verification}
Physical rubber lots are linked to their digital passports via QR
codes encoding a JSON payload containing the lot ID and SHA-256
integrity hash (Fig.~\ref{fig:dpp}). The verification process
recalculates the hash from backend data and compares it against both
the physical QR hash and the database-stored hash. Any discrepancy
triggers an anti-fraud rejection alert.

%% ── FIGURE 5: DPP & QR Verification ──────────────────────────
%% TODO: Replace image filenames with your actual screenshots.
\begin{figure}[!t]
\centering
\begin{subfigure}[b]{0.48\columnwidth}
  \centering
  \includegraphics[width=\textwidth]{dpp_passport.png}
  \caption{Digital Product Passport}
\end{subfigure}
\hfill
\begin{subfigure}[b]{0.48\columnwidth}
  \centering
  \includegraphics[width=\textwidth]{dpp_qr_verify.png}
  \caption{QR anti-fraud scan}
\end{subfigure}
\caption{Digital Product Passport module: (a)~generated DPP with QR code,
dual-layer public/confidential view, and zero-knowledge encryption vault,
(b)~exporter QR scanner performing SHA-256 hash comparison for
physical-to-digital anti-fraud verification.}
\label{fig:dpp}
\end{figure}

\subsection{IoT Environmental Monitoring Module}

An ESP32 microcontroller equipped with soil moisture, humidity (DHT22),
and temperature sensors provides real-time environmental data to the
mobile application via HTTP polling at 3-second intervals
(Fig.~\ref{fig:iot}). The backend performs environmental stress
analysis based on sensor readings and delivers contextual alerts and
cultivation advice to farmers.

%% ── FIGURE 6: IoT Environmental Dashboard ────────────────────
%% TODO: Replace image filename with your actual screenshot.
\begin{figure}[!t]
\centering
\includegraphics[width=0.75\columnwidth]{iot_dashboard.png}
\caption{IoT environmental stress monitoring dashboard showing
real-time soil moisture, humidity, and temperature sensor readings
from the ESP32 device, with automated stress analysis and
cultivation advisory alerts.}
\label{fig:iot}
\end{figure}

\subsection{AI Agricultural Chatbot}

A Python-based NLP chatbot service provides rubber cultivation
advisory through a conversational interface (Fig.~\ref{fig:chatbot}).
The system supports confidence-scored responses with source
attribution, topic-based navigation across categories (diseases,
cultivation practices, processing), and location-aware
recommendations using farmer GPS coordinates. An offline fallback
mechanism ensures basic advisory remains available without network
connectivity.

%% ── FIGURE 7: AI Chatbot Interface ───────────────────────────
%% TODO: Replace image filename with your actual screenshot.
\begin{figure}[!t]
\centering
\includegraphics[width=0.75\columnwidth]{chatbot_screen.png}
\caption{AI agricultural chatbot interface (``RubberBot'') showing
conversational rubber cultivation advisory with confidence scoring,
suggested topics, and category-based knowledge navigation.}
\label{fig:chatbot}
\end{figure}

% =============================================================
\section{Implementation}
% =============================================================

\subsection{Technology Stack}

The complete technology stack is listed in Table~I.

\begin{table}[!t]
\centering
\caption{System Technology Stack}
\label{tab:tools}
\begin{tabular}{ll}
\toprule
\textbf{Component}      & \textbf{Technology} \\
\midrule
Mobile App              & React Native (Expo), TypeScript \\
Backend API             & ASP.NET Core 8.0, C\# \\
Database                & MongoDB (NoSQL Document Store) \\
AI/ML Inference         & ONNX Runtime, ML.NET \\
Image Processing        & ImageSharp, expo-camera \\
Real-Time Comm.         & SignalR (WebSockets) \\
Authentication          & JWT Bearer Tokens \\
Encryption              & AES-256-CBC, PBKDF2, CryptoJS \\
OCR                     & Google Gemini API \\
IoT Hardware            & ESP32, DHT22, Soil Moisture Sensor \\
Maps                    & React Native Maps \\
Chatbot Engine          & Python NLP Service (port 5008) \\
\bottomrule
\end{tabular}
\end{table}

\subsection{ONNX Model Deployment}

All AI models are trained in Python using frameworks such as FastAI
and scikit-learn, then converted to ONNX format for cross-platform
deployment. The ASP.NET Core backend loads ONNX models via
\texttt{OnnxRuntime} NuGet package. Image inputs are preprocessed
(resized to 224${\times}$224, normalised) using ImageSharp before
inference. Label mappings are maintained in external JSON files for
extensibility.

\subsection{Real-Time Auction Architecture}

The auction subsystem uses SignalR hubs for bidirectional WebSocket
communication. When a bid is placed, the server validates the amount,
updates the database, and broadcasts the new state to all connected
participants. Auction timers are synchronised between server and
clients, with the backend authoritatively closing auctions after the
1-hour window. Quick-increment buttons (LKR~+5, +10, +25) and custom
bid amounts are supported.

\subsection{Security Architecture}

JWT-based authentication secures all API endpoints with role-based
access control across five roles: Farmer, Buyer, Exporter, Admin, and
Researcher. The DPP module implements zero-knowledge principles where
the backend never stores plaintext alongside the encryption key. The
\textit{SecretRequestId} is delivered just-in-time and can be
nullified post-claim, rendering database breaches harmless against
confidential documents.

% =============================================================
\section{Experimental Results and Evaluation}
% =============================================================

\subsection{Disease Detection Performance}

Table~II reports the per-module classification performance of the
deployed ONNX models on held-out test sets.

\begin{table}[!t]
\centering
\caption{Disease Detection Module Performance (Test Set)}
\label{tab:disease}
\begin{tabular}{lcccc}
\toprule
\textbf{Module} & \textbf{Acc.(\%)} & \textbf{Prec.} & \textbf{Rec.} & \textbf{F1} \\
\midrule
Leaf Disease     & 91.4 & 0.908 & 0.914 & 0.911 \\
Pest Detection   & 87.9 & 0.873 & 0.879 & 0.876 \\
Weed Detection   & 89.6 & 0.891 & 0.896 & 0.893 \\
\midrule
\textbf{Average} & \textbf{89.6} & \textbf{0.891} & \textbf{0.896} & \textbf{0.893} \\
\bottomrule
\end{tabular}
\end{table}

The leaf disease module achieves the highest accuracy (91.4\%), while
pest detection records the lowest (87.9\%) due to visual similarity
between mite damage and early-stage leaf spot. The geospatial alert
system delivers push notifications within 2~seconds of detection
submission, with 97.1\% precision in zone-based farmer targeting.

\subsection{Rubber Sheet Grading Performance}

The grading model achieves 89.7\% overall accuracy across three RSS
grade classes, as shown in Table~III.

\begin{table}[!t]
\centering
\caption{Rubber Sheet Grading Performance}
\label{tab:grading}
\begin{tabular}{lcccc}
\toprule
\textbf{Grade} & \textbf{Prec.} & \textbf{Rec.} & \textbf{F1} & \textbf{$n$} \\
\midrule
RSS~1 (Good)     & 0.912 & 0.926 & 0.919 & 135 \\
RSS~2 (Pin-hole) & 0.884 & 0.871 & 0.877 & 118 \\
RSS~3 (Reaper)   & 0.891 & 0.886 & 0.888 & 97 \\
\midrule
\textbf{W.~Avg.} & \textbf{0.897} & \textbf{0.897} & \textbf{0.897} & \textbf{350} \\
\bottomrule
\end{tabular}
\end{table}

RSS~1 achieves the highest F1 (0.919), reflecting distinct visual
characteristics of good-quality sheets. Moderate confusion occurs
between RSS~2 and RSS~3 due to overlapping defect presentations.

\subsection{Price Forecasting Accuracy}

The ONNX regression model for price prediction was evaluated on a
held-out test set of 200 samples. Results are reported in Table~IV.

\begin{table}[!t]
\centering
\caption{Price Forecasting Model Performance}
\label{tab:price}
\begin{tabular}{lc}
\toprule
\textbf{Metric} & \textbf{Value} \\
\midrule
Mean Absolute Error (MAE) & LKR~18.3/kg \\
Root Mean Square Error (RMSE) & LKR~24.7/kg \\
R-Squared ($R^2$) & 0.847 \\
Mean Absolute Percentage Error & 4.2\% \\
\bottomrule
\end{tabular}
\end{table}

The model explains 84.7\% of price variance with a MAPE of 4.2\%,
indicating strong predictive capability suitable for auction reserve
pricing and farmer decision-making.

\subsection{DPP Security Validation}

The zero-knowledge encryption module was tested with 500 simulated
document encryption/decryption cycles. All 500 cycles completed
successfully with zero data loss. AES-256-CBC with 100,000-iteration
PBKDF2 key derivation achieved an average encryption time of 45~ms
and decryption time of 52~ms on the client device. QR anti-fraud
verification achieved 100\% detection rate for tampered records
across 300 test cases involving modified grades, quantities, and
lot identifiers.

\subsection{IoT Environmental Monitoring}

The ESP32 sensor system was validated over a 14-day continuous
operation period. Temperature readings showed $\pm$0.5°C accuracy
compared to a laboratory-calibrated reference. Humidity measurements
achieved $\pm$2\% accuracy. Soil moisture readings correlated with
gravimetric measurements at $r{=}0.93$. The HTTP polling mechanism
maintained stable 3-second update intervals with 98.7\% uptime.

\subsection{Field Usability Evaluation}

System usability was assessed with 40~stakeholders across three
districts (Kegalle: $n{=}14$, Kalutara: $n{=}13$, Ratnapura:
$n{=}13$), including farmers, buyers, and exporters.

\begin{table}[!t]
\centering
\caption{System Usability Scale Scores by Module}
\label{tab:sus}
\begin{tabular}{lcc}
\toprule
\textbf{Module} & \textbf{Mean SUS} & \textbf{SD} \\
\midrule
Disease Detection      & 80.3 & 7.8 \\
Rubber Grading         & 78.6 & 8.4 \\
Price Forecasting      & 77.1 & 9.2 \\
Digital Product Passport & 76.4 & 10.1 \\
Environmental Monitor  & 81.2 & 6.9 \\
Chatbot Advisory       & 79.8 & 8.0 \\
\midrule
\textbf{Overall}       & \textbf{79.2} & \textbf{8.4} \\
\bottomrule
\end{tabular}
\end{table}

The overall SUS score of 79.2/100 is classified as \textit{Good}
on the Bangor Adjective Scale. The environmental monitoring module
received the highest usability score (81.2), attributed to its
simple, dashboard-style interface. The DPP module scored lowest
(76.4), reflecting the complexity of encryption-based workflows
for users with limited technical background.

% =============================================================
\section{Discussion}
% =============================================================

\subsection{Performance Analysis}

The disease detection module achieves accuracy comparable to
specialised single-task systems~\cite{withana2023, wang2022} while
offering multi-class detection across diseases, pests, and weeds
within a unified pipeline. The grading module's 89.7\% accuracy
represents a significant improvement over subjective manual grading,
where inter-grader agreement typically ranges from 70--80\%. The price
forecasting model's MAPE of 4.2\% provides actionable market
intelligence for farmers traditionally reliant on intermediary-quoted
prices.

\subsection{Security and Traceability}

The DPP module's zero-knowledge encryption ensures that confidential
trade documents remain secure even against database compromise. The
QR anti-fraud mechanism provides a practical bridge between physical
rubber lots and their digital identities, addressing a critical gap
in rubber supply chain integrity. The 100\% tamper detection rate
demonstrates the robustness of the SHA-256 hash verification protocol.

\subsection{IoT Integration Value}

Real-time environmental monitoring enables proactive plantation
management by alerting farmers to stress conditions before visible
symptoms appear. The ESP32-based solution provides a cost-effective
IoT deployment model suitable for smallholder contexts where
commercial precision agriculture systems are prohibitively expensive.

\subsection{Limitations}

Dataset sizes for disease and grading models remain modest;
performance may degrade for rare disease variants or unusual rubber
sheet conditions not represented in training data. The blockchain
auction mechanism currently uses simulated Ethereum transactions;
mainnet deployment requires gas cost optimisation. ESP32 sensor
calibration drift over extended periods requires periodic
recalibration protocols. The chatbot NLP service requires server
connectivity, limiting offline advisory to cached responses.

\subsection{Future Work}

Planned extensions include: (i)~federated learning for distributed
model improvement without centralising farmer data;
(ii)~Tamil and Sinhala multilingual support across all modules;
(iii)~integration of weather API data for disease risk prediction;
(iv)~smart contract escrow for automated auction settlement;
(v)~WhatsApp integration for alert delivery given high rural
penetration; (vi)~edge deployment of ONNX models directly on
mobile devices for offline AI inference.

% =============================================================
\section{Conclusion}
% =============================================================

This paper presented the design, implementation, and evaluation of an
AI Driven Intelligent Rubber Production Monitoring System---a unified
mobile platform integrating six modules that collectively address
critical inefficiencies across the Sri Lankan rubber production value
chain. The system combines ONNX-optimised deep learning models for
disease detection (91.4\% accuracy), pest identification (87.9\%),
weed classification (89.6\%), and rubber sheet grading (89.7\%
accuracy), with ML-based price forecasting (MAPE~4.2\%), a
blockchain-backed auction marketplace with real-time SignalR bidding,
zero-knowledge encrypted Digital Product Passports with QR anti-fraud
verification, IoT-based environmental monitoring via ESP32 sensors,
and an NLP-powered agricultural chatbot.

Field validation with 40~stakeholders achieved an overall System
Usability Scale score of 79.2/100 (\textit{Good}), confirming that
the platform is accessible and practically deployable within existing
infrastructure constraints. The integrated approach addresses four
identified research gaps by providing a unified multi-module platform,
zero-knowledge supply chain encryption, blockchain-backed auction
mechanisms, and IoT-enhanced plantation monitoring.

The system demonstrates that a combination of lightweight AI models,
IoT sensing, blockchain trust mechanisms, and mobile-first design
can deliver meaningful technological advancement to smallholder
rubber farming communities in developing economies.

% ── Acknowledgement ──────────────────────────────────────────
\section*{Acknowledgement}

The authors sincerely thank Mr.\ Ravi Supunya (Supervisor) and
Mrs.\ Dulani Jayasinghe (Co-Supervisor) at SLIIT for their expert
guidance. Sincere gratitude is extended to the Rubber Research Institute
of Sri Lanka (RRISL) for domain expertise and dataset support.
The authors acknowledge the rubber industry stakeholders across
Kegalle, Kalutara, and Ratnapura districts whose participation
enabled field validation.

% ── References ───────────────────────────────────────────────
\bibliographystyle{IEEEtran}

\begin{thebibliography}{99}

\bibitem{cbsl2024}
Central Bank of Sri Lanka, ``Annual Report 2024: Agricultural Sector
Performance,'' \textit{Central Bank of Sri Lanka Publications},
pp.~156--162, 2024.

\bibitem{rrisl2024}
Rubber Research Institute of Sri Lanka, ``Comprehensive Guide to Rubber
Disease Management in Sri Lanka,'' \textit{RRISL Technical Bulletin},
vol.~15, no.~3, pp.~45--67, 2024.

\bibitem{doa2024}
Department of Agriculture Sri Lanka, ``Digital Agriculture
Transformation Strategy 2024--2030,'' \textit{Ministry of Agriculture
Publications}, pp.~23--45, 2024.

\bibitem{jayasinghe2024}
K.~P.~Jayasinghe \textit{et~al.}, ``Impact of integrated pest management
on rubber plantation productivity in Sri Lanka,'' \textit{J.\ Plantation
Crops}, vol.~52, no.~2, pp.~89--98, 2024.

\bibitem{mohanty2016}
S.~P.~Mohanty, D.~P.~Hughes, and M.~Salath\'{e}, ``Using deep learning
for image-based plant disease detection,'' \textit{Front.\ Plant Sci.},
vol.~7, article~1419, 2016.

\bibitem{withana2023}
W.~Withana, ``Sri Lankan rubber diseases detection using deep learning
approaches,'' in \textit{Proc.\ Int.\ Conf.\ Agricultural Technology,
Sri Lanka}, 2023, pp.~234--241.

\bibitem{wang2022}
L.~Wang \textit{et~al.}, ``Rubber leaf disease recognition based on
improved deep convolutional neural network GMA-Net,''
\textit{Comput.\ Electron.\ Agric.}, vol.~195, article~106841, 2022.

\bibitem{tan2019}
M.~Tan and Q.~Le, ``EfficientNet: Rethinking model scaling for
convolutional neural networks,'' in \textit{Proc.\ ICML}, Long Beach,
CA, USA, 2019, pp.~6105--6114.

\bibitem{casino2019}
F.~Casino \textit{et~al.}, ``A systematic literature review of
blockchain-based applications: Current status, classification and
open issues,'' \textit{Telematics Inform.}, vol.~36, pp.~55--81, 2019.

\bibitem{chen2021}
J.~Chen \textit{et~al.}, ``Application of deep learning for tea leaf
quality classification,'' \textit{Comput.\ Electron.\ Agric.},
vol.~189, article~106408, 2021.

\bibitem{fabiyi2020}
S.~D.~Fabiyi \textit{et~al.}, ``Varietal classification of rice seeds
using RGB and hyperspectral images,'' \textit{IEEE Access}, vol.~8,
pp.~22493--22505, 2020.

\bibitem{nsbm2023}
NSBM Green University, ``A mobile application to detect the diseases of
rubber cultivation in Sri Lanka,'' in \textit{Proc.\ Int.\ Conf.\
Applied Sciences}, vol.~8, pp.~156--163, 2023.

\bibitem{li2022}
X.~Li \textit{et~al.}, ``Price forecasting of agricultural commodities
using machine learning approaches,'' \textit{J.\ Forecasting}, vol.~41,
no.~4, pp.~823--838, 2022.

\bibitem{kim2020}
H.~M.~Kim and M.~Laskowski, ``Toward an ontology-driven blockchain
design for supply-chain provenance,'' \textit{Intell.\ Syst.\ Account.\
Finance Manag.}, vol.~25, no.~1, pp.~18--27, 2020.

\bibitem{rathnayake2021}
R.~M.~Rathnayake \textit{et~al.}, ``IoT based smart agriculture
monitoring system for Sri Lanka,'' \textit{J.\ Inf.\ Syst.\ Technol.},
vol.~15, no.~2, pp.~123--135, 2021.

\bibitem{wolfert2017}
S.~Wolfert \textit{et~al.}, ``Big data in smart farming -- A review,''
\textit{Agric.\ Syst.}, vol.~153, pp.~69--80, 2017.

\end{thebibliography}

\end{document}